% Generated draft: requires mathematical and source review.
\documentclass[11pt]{article}
\usepackage[margin=25mm]{geometry}
\usepackage{fontspec}
\setmainfont{Latin Modern Roman}
\usepackage{amsmath,amssymb,mathtools,unicode-math}
\setmathfont{Latin Modern Math}
\usepackage{enumitem,needspace,xurl,hyperref}
\setlength{\parindent}{0pt}
\setlength{\parskip}{5pt}
\setlength{\emergencystretch}{3em}
\newcommand{\N}{\mathbb{N}}
\newcommand{\Z}{\mathbb{Z}}
\newcommand{\Q}{\mathbb{Q}}
\newcommand{\R}{\mathbb{R}}
\newcommand{\C}{\mathbb{C}}
\newcommand{\F}{\mathbb{F}}
\newcommand{\A}{\mathbb{A}}
\newcommand{\PP}{\mathbb{P}}
\newcommand{\E}{\mathbb{E}}
\newcommand{\eps}{\varepsilon}
\newcommand{\dd}{\,\mathrm{d}}
\begin{document}
\section*{UnsolvedMath ID: 800729}
\textbf{Source problem code:} catalogue entry
\subsection{Definitions and mathematical statement}
Let \(X\) be a topological space, and let
\[
S(X)=\{f:X\to X\mid f\text{ is continuous}\}.
\]
Equip \(S(X)\) with composition, \((f\circ g)(x)=f(g(x))\), and the subspace topology inherited from the product \(X^X=\prod_{x\in X}X\). Thus basic neighborhoods of \(f\in S(X)\) have the form
\[
\{g\in S(X):g(x_i)\in V_i\text{ for }1\leq i\leq n\},
\]
where \(n\) is finite, \(x_i\in X\), and each \(V_i\) is an open neighborhood of \(f(x_i)\). The identity element is \(\operatorname{id}_X\). A space \(X\) is homogeneous if, for every \(x_1,x_2\in X\), there is a homeomorphism \(h:X\to X\) satisfying \(h(x_1)=x_2\).

\textbf{Source question (Protasov).} Is \(S(X)\) left bounded for every zero-dimensional compact homogeneous space \(X\)?

\textbf{Unresolved conventions and conditional expansion.} The supplied record does not define left boundedness. If the intended definition for this monoid is that every neighborhood \(U\) of \(\operatorname{id}_X\) admits a finite set \(F\subseteq S(X)\) such that
\[
S(X)=F\circ U
     =\{a\circ u:a\in F,\ u\in U\},
\]
then the requested conclusion expands to
\[
\forall U\text{ a neighborhood of }\operatorname{id}_X\quad
\exists F\subseteq S(X)\text{ finite}\quad
\forall g\in S(X)\quad
\exists a\in F\ \exists u\in U\quad g=a\circ u.
\]
This is a conditional interpretation requiring confirmation from the source, rather than a supplied definition. Likewise, interpreting zero-dimensionality as the existence of a base of sets that are both open and closed requires confirmation of the source conventions. Compactness means that every open cover has a finite subcover, but the record does not specify whether a Hausdorff assumption is implicit. The intended separation assumptions and whether the empty space is admitted remain unresolved. No joint continuity of composition is asserted here.
\subsection{Short English statement}
The question attributed to Protasov asks whether the continuous selfmaps of every zero-dimensional compact homogeneous space form a left bounded semigroup under composition and pointwise convergence. The supplied registry labels the problem open. The definition of left boundedness and the source's separation conventions must be confirmed before the displayed conditional expansion can be treated as a definitive formulation.
\subsection{Sources}
\begin{itemize}
\item Local UnsolvedMath registry, permanent ID \(800729\), exact source problem code \texttt{catalogue entry}; supplied title: \emph{Open Problems in Topological Algebra --- Problem 3 (Protasov)}. The registry lists Banakh and Choban as proposers, gives the year 2012, and labels the record open. This is supplied registry information, not an independently verified status assessment.
\item Banakh and Choban, \emph{Some Open Problems in Topological Algebra} (2012), Problem 3, \url{https://arxiv.org/abs/1202.1619}. The supplied background attributes the problem to Protasov and identifies the extraction method as \texttt{arxiv\_latex}.
\item The supplied background records an access date of 2026-10-16. This is supplied provenance; no independent source visit or verification is claimed.
\end{itemize}
\end{document}
